% !TeX root = ../main.tex
\chapter{Funzioni: definizioni e proprietà generali}
\label{cap:funzioni}
\fonti{Lezioni 5 e 6}

\section{Dominio, codominio e immagine}
% Fonte: Lezione 5, p.~12

\begin{definizione}
Una \emph{funzione} $f\colon X\to Y$ è una legge che associa a ogni
elemento $x\in X$ uno e un solo elemento $y\in Y$. Scriviamo
\[
f\colon X\longrightarrow Y,\qquad x\longmapsto y=f(x).
\]
L'insieme $X$ è il \emph{dominio}, $Y$ il \emph{codominio}; $f(x)$ è
l'\emph{immagine di $x$ tramite $f$}. L'\emph{insieme immagine} è
\[
\operatorname{Im}(f)=f(X)=\{f(x):x\in X\}\subseteq Y.
\]
\end{definizione}

La condizione «uno e un solo» riguarda il valore associato a ciascun
elemento del dominio. Elementi distinti possono avere la stessa immagine;
inoltre un elemento del codominio può non essere raggiunto. Perciò
codominio e immagine non coincidono necessariamente.

\section{Iniettività, suriettività e biettività}
% Fonte: Lezione 5, p.~13

\begin{definizione}
Una funzione $f\colon X\to Y$ si dice:
\begin{itemize}
  \item \emph{iniettiva} se elementi distinti hanno immagini distinte:
  \[
  \forall x_1,x_2\in X,\quad
  x_1\ne x_2\ \Longrightarrow\ f(x_1)\ne f(x_2);
  \]
  equivalentemente, $f(x_1)=f(x_2)$ implica $x_1=x_2$;
  \item \emph{suriettiva} se ogni elemento del codominio è raggiunto:
  \[
  \forall y\in Y\ \exists x\in X:\ f(x)=y,
  \qquad\text{ossia }\operatorname{Im}(f)=Y;
  \]
  \item \emph{biettiva} se è sia iniettiva sia suriettiva.
\end{itemize}
\end{definizione}

Per una biezione, dunque, l'equazione $f(x)=y$ ha esattamente una
soluzione $x\in X$ per ogni $y\in Y$: la suriettività dà l'esistenza,
l'iniettività dà l'unicità.

\begin{esempio}
La funzione $f\colon\R\to\R$, $f(x)=x^2$, non è iniettiva perché
$f(1)=f(-1)=1$, e non è suriettiva perché non assume valori negativi.
Se il codominio diventa $[0,+\infty)$, la funzione è suriettiva.
Restringendo anche il dominio a $[0,+\infty)$ diventa biettiva.
Dominio e codominio sono quindi parte essenziale della descrizione di
una funzione.
\end{esempio}

\section{Composizione e suo dominio}
% Fonte: Lezione 5, pp.~14--15

Siano $f\colon X\to Y$ e $g\colon V\to Z$. Per calcolare $g(f(x))$
occorre che il risultato della prima funzione appartenga al dominio
della seconda. Poniamo pertanto
\[
X'=\{x\in X:f(x)\in V\}.
\]
La \emph{funzione composta} è
\[
g\circ f\colon X'\longrightarrow Z,\qquad
(g\circ f)(x)=g(f(x)).
\]
Se $f(X)\subseteq V$, il dominio della composta è tutto $X$; se c'è
soltanto un'intersezione non vuota fra $f(X)$ e $V$, la composta è
definita sul sottoinsieme $X'$. In assenza di intersezione, $X'$ è vuoto.
L'ordine di lettura è da destra verso sinistra: si applica prima $f$ e
poi $g$.

\begin{schema}
\begin{tikzpicture}[>=Stealth]
  \node (X) at (0,0) {$X'$};
  \node (V) at (4,0) {$V$};
  \node (Z) at (8,0) {$Z$};
  \draw[->,coloreschema] (X) -- node[above] {$f|_{X'}$} (V);
  \draw[->,coloreschema] (V) -- node[above] {$g$} (Z);
  \draw[->] (X) to[bend right=22] node[below] {$g\circ f$} (Z);
\end{tikzpicture}
\end{schema}

\begin{esempio}
Prendiamo $f(x)=1+x^2$, con dominio $\R$, e $g(y)=\sqrt y$, con
dominio $[0,+\infty)$. Si ottiene
\[
(f\circ g)(x)=1+(\sqrt x)^2=1+x,
\qquad \operatorname{dom}(f\circ g)=[0,+\infty),
\]
mentre
\[
(g\circ f)(x)=\sqrt{1+x^2},
\qquad \operatorname{dom}(g\circ f)=\R.
\]
La composizione, in generale, non è commutativa: qui differiscono sia
le formule sia i domini. Anche quando una formula si semplifica, il
dominio della composta resta quello imposto dalle operazioni originarie.
\end{esempio}

\section{Grafico, controimmagini e funzione inversa}
% Fonte: Lezione 5, p.~16; lezione 6, p.~2

\begin{definizione}
Il \emph{grafico} di $f\colon X\to Y$ è l'insieme
\[
\operatorname{graf}(f)=\{(x,f(x)):x\in X\}\subseteq X\times Y.
\]
La \emph{controimmagine}, o \emph{fibra}, di $y\in Y$ è
\[
f^{-1}(\{y\})=\{x\in X:f(x)=y\}.
\]
Più in generale, per $B\subseteq Y$ si pone
$f^{-1}(B)=\{x\in X:f(x)\in B\}$.
\end{definizione}

La controimmagine è un insieme ed è definita per ogni funzione. Se $f$
è iniettiva, la fibra di un elemento $y\in\operatorname{Im}(f)$ contiene
esattamente un elemento: non è vuota perché $y$ appartiene all'immagine,
e non ne può contenere due per l'iniettività. Questo consente di
definire una nuova funzione.

\begin{definizione}
Se $f\colon X\to Y$ è iniettiva, la \emph{funzione inversa} è
\[
f^{-1}\colon\operatorname{Im}(f)\longrightarrow X,
\qquad f^{-1}(y)=\text{l'unico }x\in X\text{ tale che }f(x)=y.
\]
\end{definizione}

La scrittura $f^{-1}(y)$ indica qui un elemento, mentre
$f^{-1}(\{y\})=\{f^{-1}(y)\}$ indica il corrispondente insieme con un
solo elemento. Inoltre $f^{-1}$ non indica il reciproco $1/f$.
L'inversa è definita sull'intero codominio $Y$ se $f$ è biettiva.

\begin{proposizione}
L'inversa di una funzione iniettiva soddisfa
\[
f^{-1}\circ f=\operatorname{Id}_X,
\qquad f\circ f^{-1}=\operatorname{Id}_{\operatorname{Im}(f)},
\]
dove $\operatorname{Id}_A\colon A\to A$, $a\mapsto a$, è la funzione
identità su $A$.
\end{proposizione}
\begin{proof}
Per $x\in X$, l'unico elemento che $f$ manda in $f(x)$ è $x$ stesso;
quindi $f^{-1}(f(x))=x$. Per $y\in\operatorname{Im}(f)$, la definizione
di $f^{-1}(y)$ dà direttamente $f(f^{-1}(y))=y$.
\end{proof}

Quando $X,Y\subseteq\R$, il punto $(x,y)$ appartiene al grafico di $f$
se e solo se $(y,x)$ appartiene al grafico dell'inversa: i due grafici
sono simmetrici rispetto alla bisettrice $y=x$.

\begin{schema}
\begin{tikzpicture}[x=1cm,y=1cm]
  \draw[asse] (-.3,0) -- (4.7,0) node[right] {$x$};
  \draw[asse] (0,-.3) -- (0,4.7) node[above] {$y$};
  \draw[dashed,black!45] (0,0) -- (4.4,4.4) node[above right] {$y=x$};
  \draw[intervallo,domain=0:2.05,samples=60] plot (\x,{\x*\x});
  \draw[thick,orange!80!black,domain=0:2.05,samples=60]
    plot ({\x*\x},\x);
  \node[coloreschema,left] at (1.8,3.7) {$y=x^2,\ x\ge0$};
  \node[orange!80!black,below] at (3.4,1.7) {$y=\sqrt x$};
  \fill[coloreschema] (1.5,2.25) circle (1.5pt);
  \fill[orange!80!black] (2.25,1.5) circle (1.5pt);
  \draw[dotted] (1.5,2.25) -- (2.25,1.5);
\end{tikzpicture}
\end{schema}

% !TeX root = ../main.tex
\section{Funzioni reali e dominio naturale}
% Fonte: Lezione 6, p.~3

Una \emph{funzione reale di variabile reale} ha dominio e codominio contenuti in $\R$: $f\colon X\to Y$, con $X,Y\subseteq\R$. Quando è assegnata mediante una formula, il suo \emph{dominio naturale} è il più grande sottoinsieme di $\R$ in cui tutte le operazioni della formula hanno significato reale.

\begin{esempio}
Per $f(x)=1/(x^2-4)$ occorre $x^2-4\ne0$. Dunque
\[
\operatorname{dom}(f)=\R\setminus\{-2,2\}.
\]
I punti esclusi dal dominio non tornano a farne parte per una semplice manipolazione algebrica della formula.
\end{esempio}

\section{Operazioni tra funzioni}
% Fonte: Lezione 6, p.~4

Date $f,g\colon D\subseteq\R\to\R$, le operazioni sono definite \emph{punto per punto}:
\[
\begin{aligned}
(f+g)(x)&=f(x)+g(x), & (f-g)(x)&=f(x)-g(x),\\
(fg)(x)&=f(x)g(x), & \left(\frac fg\right)(x)&=\frac{f(x)}{g(x)}.
\end{aligned}
\]
\textbf{Domini delle operazioni.}
\begin{itemize}
  \item Somma, differenza e prodotto hanno dominio $D$.
  \item Il quoziente è definito in $\{x\in D:g(x)\ne0\}$.
  \item Se i domini iniziali sono diversi, si usa la loro intersezione; per il quoziente si escludono inoltre gli zeri del denominatore.
\end{itemize}
Il prodotto $fg$ e la composizione $f\circ g$ sono operazioni diverse.

\section{Monotonia e composizione}
% Fonte: Lezione 6, pp.~5--7

\begin{definizione}
Sia $f\colon D\subseteq\R\to\R$.
\begin{itemize}
  \item È \emph{crescente}, o \emph{non decrescente}, se
\[
\forall x_1,x_2\in D,\quad x_1<x_2\ \Longrightarrow\ f(x_1)\le f(x_2).
\]
È \emph{strettamente crescente} se, sotto la stessa ipotesi, $f(x_1)<f(x_2)$.
  \item È \emph{decrescente}, o \emph{non crescente}, se
\[
x_1<x_2\ \Longrightarrow\ f(x_1)\ge f(x_2),
\]
ed è \emph{strettamente decrescente} se vale $f(x_1)>f(x_2)$.
\end{itemize}
Si dice \emph{monotona} se è crescente oppure decrescente.
\end{definizione}

\textbf{Confronto sull’intero dominio.} La monotonia confronta tutte le coppie ordinate di punti del dominio. Una funzione costante è sia crescente sia decrescente nel senso non stretto; su un dominio con almeno due punti non è strettamente monotona.

\begin{schema}
\begin{tikzpicture}[x=1.5cm,y=.85cm]
  \draw[asse] (-.2,0) -- (4.6,0) node[right] {$x$};
  \draw[asse] (0,-.2) -- (0,3.7) node[above] {$y$};
  \draw[intervallo] (.2,.4) -- (4.1,2.9);
  \draw[thick,orange!80!black] (.2,3.2) -- (4.1,.6);
  \draw[thick,black!55] (.2,3.4) -- (4.1,3.4);
  \node[coloreschema,right] at (4.1,2.9) {crescente};
  \node[orange!80!black,right] at (4.1,.6) {decrescente};
  \node[black!65,right] at (4.1,3.4) {costante};
\end{tikzpicture}
\end{schema}

\begin{proposizione}
\begin{enumerate}[label=\arabic*.]
\item La composizione di due funzioni monotone è monotona: se hanno lo stesso verso è crescente, se hanno versi opposti è decrescente.
\item Se entrambe sono strettamente monotone, anche la composta è strettamente monotona.
\end{enumerate}
\end{proposizione}
\begin{proof}
Scriviamo la composta come $g\circ f$ e scegliamo $x_1<x_2$ nel suo dominio.

\begin{enumerate}[label=\arabic*.]
\item \textbf{Stesso verso.} Se $f$ e $g$ sono crescenti, da $f(x_1)\le f(x_2)$ segue $g(f(x_1))\le g(f(x_2))$. Se entrambe sono decrescenti, si ha $f(x_1)\ge f(x_2)$; applicando $g$ si inverte nuovamente l'ordine e si ottiene ancora $g(f(x_1))\le g(f(x_2))$.

\textbf{Versi opposti.} Se $f$ è crescente e $g$ decrescente, da $f(x_1)\le f(x_2)$ segue $g(f(x_1))\ge g(f(x_2))$. Se $f$ è decrescente e $g$ crescente, da $f(x_1)\ge f(x_2)$ si ottiene la stessa conclusione. La composta è quindi decrescente.
\item \textbf{Caso stretto.} Quando entrambe le funzioni sono strettamente monotone, tutte le disuguaglianze corrispondenti sono strette.
\end{enumerate}
\end{proof}

\section{Stretta monotonia e invertibilità}
% Fonte: Lezione 6, pp.~7--8

\begin{proposizione}
\begin{enumerate}[label=\arabic*.]
\item Ogni funzione strettamente monotona è iniettiva.
\item Se una funzione monotona è iniettiva, è strettamente monotona e la sua inversa sull'immagine ha lo stesso verso di monotonia.
\end{enumerate}
\end{proposizione}
\begin{proof}
\begin{enumerate}[label=\arabic*.]
\item \textbf{Iniettività.} Per $x_1\ne x_2$, uno dei due punti è minore dell'altro. La stretta monotonia impone valori distinti: questo prova l'iniettività.
\item \textbf{Stretta monotonia.} In una funzione monotona iniettiva, l'uguaglianza fra due valori in punti distinti è esclusa, quindi le disuguaglianze diventano strette.

\textbf{Inversa crescente.} Supponiamo $f$ strettamente crescente e siano $y_1<y_2$ nella sua immagine. Posti $x_i=f^{-1}(y_i)$, non può essere $x_1\ge x_2$: si avrebbe $y_1=f(x_1)\ge f(x_2)=y_2$. Dunque $x_1<x_2$ e l'inversa è strettamente crescente.

\textbf{Inversa decrescente.} Se $f$ è strettamente decrescente e $y_1<y_2$, non può essere $x_1\le x_2$: si avrebbe $y_1=f(x_1)\ge f(x_2)=y_2$. Dunque $x_1>x_2$ e l'inversa è strettamente decrescente.
\end{enumerate}
\end{proof}

\begin{osservazione}
L'iniettività da sola non impone la monotonia. Per $f(x)=1/x$ su $\R\setminus\{0\}$:
\begin{itemize}
  \item \textbf{Iniettività.} Da $1/x_1=1/x_2$ segue $x_1=x_2$.
  \item \textbf{Decrescenza sulle semirette.} La funzione è decrescente su ciascuna delle due semirette separatamente; non lo è sull'intero dominio, perché $-1<1$ ma $f(-1)=-1<1=f(1)$.
  \item \textbf{Mancanza di crescenza.} Non è neppure crescente: $1<2$ ma $f(1)>f(2)$.
\end{itemize}
\end{osservazione}

\begin{schema}
\begin{tikzpicture}[scale=.85]
  \draw[asse] (-3.5,0) -- (3.5,0) node[right] {$x$};
  \draw[asse] (0,-3.1) -- (0,3.2) node[above] {$y$};
  \draw[intervallo,domain=.35:3.2,samples=80] plot (\x,{1/\x});
  \draw[intervallo,domain=-3.2:-.35,samples=80] plot (\x,{1/\x});
  \node[coloreschema,right] at (1.2,1.4) {$y=1/x$};
\end{tikzpicture}
\end{schema}

\section{Simmetrie: funzioni pari e dispari}
% Fonte: Lezione 6, p.~9

\textbf{Ipotesi sul dominio.} Per parlare di parità o disparità richiediamo che il dominio sia \emph{simmetrico rispetto a zero}: $x\in D$ implica $-x\in D$.

\begin{definizione}
La funzione $f\colon D\to\R$ è:
\begin{itemize}
  \item \emph{pari} se $f(-x)=f(x)$ per ogni $x\in D$;
  \item \emph{dispari} se $f(-x)=-f(x)$ per ogni $x\in D$.
\end{itemize}
\end{definizione}

\textbf{Interpretazione geometrica.}
\begin{itemize}
  \item Per una funzione pari, da $(x,f(x))$ sul grafico si ottiene $(-x,f(x))$: la simmetria è rispetto all'asse $y$.
  \item Per una funzione dispari si ottiene $(-x,-f(x))$: la simmetria è rispetto all'origine. Se $0\in D$, necessariamente $f(0)=0$, perché $f(0)=-f(0)$.
\end{itemize}
\begin{esempio}
Le funzioni $x^2$ e $x^3$ sono rispettivamente pari e dispari; i grafici sono riportati nel capitolo sulle funzioni elementari.
\end{esempio}

\section{Periodicità}
% Fonte: Lezione 6, p.~10

\begin{definizione}
Una funzione non costante $f\colon D\subseteq\R\to\R$ si dice \emph{periodica}\footnote{Si escludono qui le funzioni costanti. Per una funzione costante su $\R$, ogni $T>0$ sarebbe un periodo, senza un minimo positivo.} se esiste $T>0$ tale che, per ogni $x\in D$,
\[
x+T\in D,\quad x-T\in D,\qquad f(x+T)=f(x).
\]
Ogni tale $T$ è un \emph{periodo}. Se esiste il più piccolo periodo positivo, esso si chiama \emph{periodo principale} o fondamentale.
\end{definizione}

\textbf{Conseguenze della periodicità.}
\begin{itemize}
  \item \textbf{Traslazioni del dominio e dei valori.} Applicando la proprietà più volte, $x+kT\in D$ e $f(x+kT)=f(x)$ per ogni $k\in\Z$. Ogni multiplo intero positivo di un periodo è quindi ancora un periodo.
  \item \textbf{Mancanza di iniettività.} Una funzione periodica non costante non è iniettiva: $x$ e $x+T$ sono punti distinti con la stessa immagine.
\end{itemize}
\begin{esempio}
Seno e coseno hanno periodo principale $2\pi$; la tangente ha periodo principale $\pi$.
\end{esempio}

\section{Limitatezza ed estremi di una funzione}
% Fonte: Lezione 6, pp.~11--13

Per parlare degli estremi consideriamo un dominio $D$ non vuoto.

\begin{definizione}
Una funzione $f\colon D\to\R$ è \emph{limitata} se il suo insieme immagine è limitato; equivalentemente, esistono $m,M\in\R$ tali che
\[
m\le f(x)\le M\qquad\forall x\in D.
\]
Si definiscono gli estremi attraverso l'immagine:
\[
\sup_D f=\sup f(D),\qquad \inf_D f=\inf f(D).
\]
Quando esistono, $\max_D f=\max f(D)$ e $\min_D f=\min f(D)$.
\end{definizione}

\textbf{Valori e punti estremanti.}
\begin{itemize}
  \item $M=\max_D f$ è un \emph{valore massimo}; ogni $x_M\in D$ con $f(x_M)=M$ è un \emph{punto di massimo}.
  \item $m=\min_D f$ è un \emph{valore minimo}; ogni $x_m\in D$ con $f(x_m)=m$ è un \emph{punto di minimo}.
\end{itemize}
Sono estremi \emph{assoluti}: il confronto riguarda tutti i valori sul dominio. Il valore massimo o minimo, se esiste, è unico; può però essere raggiunto in più punti, anche agli estremi del dominio.

\begin{schema}
\begin{tikzpicture}[x=1.65cm,y=.85cm]
  \draw[asse] (-.2,0) -- (6,0) node[right] {$x$};
  \draw[asse] (0,-1.6) -- (0,2.9) node[above] {$y$};
  \draw[dashed,black!45] (0,2.3) node[left] {$M$} -- (5.7,2.3);
  \draw[dashed,black!45] (0,-1.1) node[left] {$m$} -- (5.7,-1.1);
  \draw[intervallo] (.8,2.3) .. controls (1.1,1.2) and (2,1.2) .. (3,2.3)
    .. controls (3.6,2.3) and (3.9,-1.1) .. (4.5,-1.1)
    .. controls (4.9,-1.1) and (5.3,-.5) .. (5.6,0);
  \foreach \x/\y in {.8/2.3,3/2.3,4.5/-1.1,5.6/0}
    \fill[coloreschema] (\x,\y) circle (1.8pt);
  \draw[dotted] (.8,0) -- (.8,2.3);
  \draw[dotted] (3,0) -- (3,2.3);
  \draw[dotted] (4.5,0) -- (4.5,-1.1);
  \node[below] at (.8,0) {$a=x_M^{(1)}$};
  \node[below] at (3,0) {$x_M^{(2)}$};
  \node[above] at (4.5,0) {$x_m$};
  \node[below] at (5.6,0) {$b$};
\end{tikzpicture}

Grafico schematico su $[a,b]$: due punti di massimo e un punto di minimo; ogni valore resta nella fascia $m\le y\le M$.
\end{schema}

\begin{esempio}
La funzione $f(x)=x$ su $[0,1)$ è limitata, ha minimo $0$ nel punto $x=0$ e supremo $1$, ma non ha massimo. La limitatezza non garantisce quindi che entrambi gli estremi siano raggiunti.
\end{esempio}

